% Reviewed solution for CritPt Challenge 56. Compiles with the supplied main.tex (natbib) and, through the
% fallbacks below, also with the original main.tex of the challenge package. Needs refs.bib, figs/ and template.py.
\providecommand{\citep}[2][]{\cite{#2}}
\providecommand{\citet}[2][]{\cite{#2}}
\providecommand{\citealt}[2][]{\cite{#2}}
\providecommand{\nat}{Nature}
\providecommand{\prl}{Physical Review Letters}
\providecommand{\prd}{Physical Review D}
\providecommand{\pra}{Physical Review A}

{\color{blue}\textbf{Reviewer's note.} The AI-generated solution has been replaced in full by the solution below, which was derived independently before the AI text was read, as the evaluation instructions suggest. The final answer follows the solution. An item-by-item assessment of the challenge statement, of the AI-generated solution and of its numerical script, with the AI text and script reproduced verbatim, follows the final answer in the accompanying document.}

\subsubsection*{Solution}

\paragraph*{Overview, conventions, and reproducibility}
The solution is provided in the following steps:
%
\begin{enumerate}[label=Step~\arabic*:,leftmargin=*,align=left,labelsep=0.5em,parsep=0pt,itemsep=0.1em,topsep=0.2em]
    \item The dark-matter field amplitude is obtained from the stated density.
    \item The photon coupling turns the field into an apparent oscillation of $\alpha$.
    \item The oscillating $\alpha$ changes the thickness and the refractive index of the beamsplitter.
    \item Those changes produce a differential optical path in the Michelson interferometer.
    \item The differential path, sampled at the finite light-travel times of a 40~km instrument, is converted to the equivalent strain against which the noise curve is quoted.
    \item The finite coherence of the field fixes how the signal-to-noise ratio (SNR) grows with observation time.
    \item The smallest detectable $\Lambda_\gamma^{-1}$ follows.
\end{enumerate}
%
Throughout, $L=40\,\mathrm{km}$ is the arm length, $l=6\,\mathrm{cm}$ and $n=3.5$ the thickness and refractive index of the beamsplitter, $\sqrt{S_h}$ the strain amplitude spectral density of the detector, $T$ the observation time, and $f=\omega/2\pi=200\,\mathrm{Hz}$ the frequency probed.

Three conventions are fixed up front so that no ambiguity remains. (i) $\Lambda_\gamma^{-1}$ is reported in $\mathrm{GeV}^{-1}$, the unit used in the interferometer literature \citep{GroteStadnik2019,Vermeulen2021}. (ii) The strain sensitivity $\sqrt{S_h}=2\times10^{-25}\,\mathrm{Hz}^{-1/2}$ is a one-sided amplitude spectral density, as all published strain-noise curves are; the SNR of a signal of amplitude $h_0$ observed for a time $T$ is then $\mathrm{SNR}^2=4\int_0^\infty|\tilde h(f)|^2/S_h(f)\,\mathrm{d}f=h_0^2T/S_h$ for a monochromatic signal (derived in Step~6). (iii) $0.7\,\mathrm{yr}$ is $0.7$ Julian years, $2.2090\times10^{7}\,\mathrm{s}$.\footnote{The choice between a Julian year ($365.25$~d) and a $365$~d year changes the final answer in the fourth significant digit only.}

Every number below is produced by \texttt{challenge56.py}, which lists all problem inputs and all model inputs (with their sources) at its top, uses \texttt{astropy} (CODATA~2018) for physical constants, ray-traces the interferometer in 40-digit arithmetic, and prints its internal consistency checks under a \texttt{--verbose} flag. The problem's own numbers are used exactly as given, even where they differ slightly from measured literature values (for example $n=3.5$ versus the measured $n=3.45$ of silicon at $2\,\mu$m \citep{Frey2006}); measured material data enter only where the problem prompt does not already provide a value explicitly.

\paragraph*{Step 1: The dark-matter field}
A coherently oscillating scalar field of mass $m_\phi$ is
%
\begin{equation}
\phi(t,\mathbf r)=\phi_0\cos(\omega t-\mathbf k\cdot\mathbf r),\qquad \omega=\frac{m_\phi c^2}{\hbar},\qquad \mathbf k=\frac{m_\phi\mathbf v}{\hbar},
\label{eq:field}
\end{equation}
%
where $\phi_0$ is the field amplitude, $\omega$ the angular Compton frequency, $\mathbf k$ the wave vector, $\mathbf v$ the velocity of the dark matter relative to the detector, $\hbar$ the reduced Planck constant, and $c$ the speed of light. At $f=\omega/2\pi=200\,\mathrm{Hz}$ as asked in this problem, we have $m_\phi c^2=hf=8.2713\times10^{-13}\,\mathrm{eV}$.\footnote{The observed frequency is $f=(m_\phi c^2+\tfrac12 m_\phi v^2)/h$; the kinetic term shifts $m_\phi$ by $\tfrac12(v/c)^2=3\times10^{-7}$, far below three significant digits and ignored.} For a monochromatic field the energy density $\rho=\tfrac12\dot\phi^2+\tfrac12 m_\phi^2\phi^2=\tfrac12 m_\phi^2\phi_0^2$ exactly in natural units, so
\begin{equation}
\phi_0=\frac{\sqrt{2\rho}}{m_\phi}.
\label{eq:phi0}
\end{equation}
With $\rho=178\times0.4\,\mathrm{GeV\,cm^{-3}}=71.2\,\mathrm{GeV\,cm^{-3}}=5.4707\times10^{-40}\,\mathrm{GeV}^4$ (converted with $(\hbar c)^3$), Eq.~\eqref{eq:phi0} yields $\phi_0=39.991\,\mathrm{GeV}$.\footnote{Because an ultralight bosonic dark-matter field such as $\phi$ is a superposition of many waves with random phases, drawn from the velocity distribution of the halo, its amplitude in any one coherence patch is Rayleigh distributed and Eq.~\eqref{eq:phi0} gives the root-mean-square value \citep{Centers2021}. ``SNR~$=1$'' below is therefore an expectation-value statement; an exclusion limit at a stated confidence level would carry an additional factor of about $3$, and up to $10$ depending on the statistical treatment, in the short-observation regime \citep{Centers2021}.}\footnote{Note that the overdensity factor $178\simeq18\pi^2$ is conventionally adopted in the astrophysics and cosmology literature as the virial overdensity of a collapsed halo relative to the cosmic mean density, and a dark-matter overdensity gravitationally bound to the Earth or the Sun would have characteristic velocities, both its internal dispersion and its streaming velocity relative to the detector, far below $230\,\mathrm{km\,s^{-1}}$ \citep[Sec.~III~F]{GroteStadnik2019}. Whether the problem's $230\,\mathrm{km\,s^{-1}}$ is meant as a dispersion or as a streaming velocity is itself not specified; the consequences are quantified in Step~6. The stated inputs do not describe a self-consistent astrophysical scenario, but they are unambiguous as inputs and are used as given.}

\paragraph*{Step 2: The coupling as an oscillating fine-structure constant}
Adding the dark-sector interaction $\mathcal L_{\rm int}=\phi F^{\mu\nu}F_{\mu\nu}/(4\Lambda_\gamma)$ to the Maxwell term $-\tfrac14F^{\mu\nu}F_{\mu\nu}$ gives the total photon kinetic term $-\tfrac14(1-\phi/\Lambda_\gamma)F^{\mu\nu}F_{\mu\nu}$. Because the coefficient of $A_\mu J^\mu$ is the physical charge only when the photon kinetic term has its canonical form $-\tfrac14F^{\mu\nu}F_{\mu\nu}$, we restore that form by the field redefinition $A_\mu\to A_\mu(1-\phi/\Lambda_\gamma)^{-1/2}$, which transfers the factor to the charge, $e\to e\,(1-\phi/\Lambda_\gamma)^{-1/2}$, so that $\alpha\to\alpha/(1-\phi/\Lambda_\gamma)$ and, to first order in $\phi/\Lambda_\gamma$ \citep[see also e.g.,][]{DamourDonoghue2010,GroteStadnik2019},
%
\begin{equation}
\epsilon(t,\mathbf r)\equiv\frac{\delta\alpha}{\alpha}(t,\mathbf r)=\frac{\phi(t,\mathbf r)}{\Lambda_\gamma}=\epsilon_0\cos(\omega t-\mathbf k\cdot\mathbf r),\qquad \epsilon_0\equiv\frac{\phi_0}{\Lambda_\gamma}=39.991\,\mathrm{GeV}\times\Lambda_\gamma^{-1}.
\label{eq:dalpha}
\end{equation}
%
Only $\alpha$ varies in Eq.~\eqref{eq:dalpha}. At tree level this is a direct consequence of the problem's Lagrangian, which couples $\phi$ to only $F^{\mu\nu}F_{\mu\nu}$; the electron mass $m_e$, which sets atomic sizes together with $\alpha$ (Step~3), is untouched. Beyond tree level, holding $m_e$ fixed is a convention rather than a consequence of the Lagrangian: a varying $\alpha$ feeds into $m_e$ through the electromagnetic self-energy, $\delta m_e/m_e\simeq(3\alpha/2\pi)\ln(\mu/m_e)\,\delta\alpha/\alpha$ with $\mu$ the scale at which the mass parameter is held fixed, a coefficient of order $10^{-2}$ whose value depends on the renormalization scheme. It is conventional instead to treat $\Lambda_\gamma$ and the electron coupling $\Lambda_e$ as independent parameters and constrain one while setting the other to zero \citep{GroteStadnik2019,Vermeulen2021}; we adopt the same definition of the photon-only scenario, $\delta m_e=0$, and carry the induced shift as a third-digit uncertainty alongside the relativistic correction of Step~3. The ultraviolet origin of the coupling (e.g., Higgs portal, dilaton, relaxion) does not enter the calculation; $\Lambda_\gamma^{-1}$ is the single parameter being bounded, and nothing gravitational is being probed.\footnote{The gravitational effect of the overdensity itself is negligible: the oscillating-metric strain of a scalar field is $h\sim G\rho/\omega^2\approx5\times10^{-36}$ at $200\,\mathrm{Hz}$, nine orders of magnitude below the detection threshold of Step~7.}

\paragraph*{Step 3: Response of the beamsplitter material}
\emph{Thickness.} In nonrelativistic quantum mechanics with Coulomb interactions, every length is proportional to the Bohr radius $a_0=\hbar/(m_e c\,\alpha)$. The thickness of a solid, $N$ lattice spacings, therefore scales exactly as $l\propto1/\alpha$. The linear relation $\delta l/l=-\delta\alpha/\alpha$ is thus exact to first order in $\delta\alpha/\alpha\sim10^{-22}$ \citep{StadnikFlambaum2015,GroteStadnik2019}. We write
\begin{equation}
\frac{\delta l}{l}=\kappa_l\frac{\delta\alpha}{\alpha},\qquad \kappa_l=-(1+K_\alpha),
\label{eq:kappal}
\end{equation}
where $K_\alpha$ collects the relativistic correction to the $\alpha$-scaling of atomic sizes. The single-particle estimate of \citet[Eq.~12]{StadnikFlambaum2015} is $K_\alpha=2(Z\alpha)^2/[\nu(j+\tfrac12)]$, with $Z$ the nuclear charge, $\nu$ the effective principal quantum number of the valence electron, and $j$ its total angular momentum. For silicon, $Z=14$, $(Z\alpha)^2=0.0104$, $\nu=\sqrt{\mathrm{Ry}/8.15\,\mathrm{eV}}=1.29$ with $\mathrm{Ry}=13.606\,\mathrm{eV}$ the Rydberg energy and $8.15\,\mathrm{eV}$ the first ionization energy, and $j=\tfrac12$ or $\tfrac32$ for the $3p$ valence shell, giving $K_\alpha=0.8\%$ to $1.6\%$, which is not negligible at the precision requested and which only a many-body calculation of the silicon lattice constant could pin down.\footnote{For lighter materials such as silica and sapphire this correction is negligible, as noted by \citet{GroteStadnik2019}.} The baseline below uses $\kappa_l=-1$, the implicit assumption of the problem, and the band $K_\alpha\in[0.008,0.016]$ is carried as a systematic.\footnote{Finite-nuclear-mass corrections affect only the $m_e/m_N$ dependence and are absent here; QED corrections to the atomic size are $\mathcal O(\alpha/\pi)$ without a large logarithm and are below the third digit. The radiative shift of $m_e$ itself (Step~2), which would enter $\delta l/l$ with unit weight, is of order $10^{-2}$ and depends on the renormalization scheme, that is, on the scale at which the electron mass parameter is held fixed (Step~2); it is set to zero by the photon-only convention and belongs with $K_\alpha$ among the third-digit uncertainties.}

Importantly, this calculated response is adiabatic. The fundamental longitudinal mode of a $6$~cm silicon plate is $f_0=v_s/2l\approx70\,\mathrm{kHz}$, with $v_s\approx8.4\,\mathrm{km\,s^{-1}}$ the longitudinal sound speed, so the driven-oscillator correction $[1-(f/f_0)^2]^{-1}$ of \citet[Eq.~7]{GroteStadnik2019} is $1+8\times10^{-6}$. The freely suspended plate expands about its center of mass, so each face moves by $\delta l/2$ along the plate normal.\footnote{Suspension resonances (pendulum and violin modes) respond to forces on the center of mass, not to a size change of the suspended optic \citep{Geraci2019,GroteStadnik2019}; they therefore do not enter the thickness term, and the gradient-force term in Step~5 is treated in the free-mass limit, valid well above the pendulum frequency.}

\emph{Refractive index.} In the same nonrelativistic scaling, all electronic excitation frequencies scale as $\omega_{\rm el}\propto m_e\alpha^2$ and the static dielectric response of the medium in units of $a_0$ is $\alpha$-independent, so the index at a fixed laser frequency $\omega_L$ is a function $n=F(\omega_L/\omega_{\rm el})$ with $F$ independent of $\alpha$ \citep[Eqs.~9 to 13]{GroteStadnik2019}. Hence
\begin{equation}
\frac{\delta n}{n}=\kappa_n\frac{\delta\alpha}{\alpha},\qquad \kappa_n=-2\,\frac{\omega_L}{n}\frac{\mathrm{d}n}{\mathrm{d}\omega_L}=+2\,\frac{\lambda_L}{n}\frac{\mathrm{d}n}{\mathrm{d}\lambda_L},
\label{eq:kappan}
\end{equation}
with $\lambda_L$ the laser wavelength. The laser frequency itself is locked to the common-mode arm length, which is isolated from size changes by the multiple-pendulum suspensions \citep{GroteStadnik2019}. The residual common-mode change of the arm length from the mirror faces is $w\,\epsilon/L\sim10^{-5}\epsilon$, with $w$ the mirror thickness, and thus negligible.\footnote{This is the point made by \citet{StadnikFlambaum2015}: if the laser were locked to a rigid resonator whose length scaled with $a_0$, adiabatic changes of $\alpha$ would be a change of units and unobservable. Gravitational-wave detectors avoid this because the laser is locked to the arm length between suspended mirrors, which does not scale with $a_0$.} The index term requires a material and a laser wavelength, neither of which the problem specifies. Among substrate materials used in gravitational-wave detectors only silicon has $n\approx3.5$ in the near-infrared ($3.45$ at $2\,\mu$m, \citealt{Frey2006}). The laser wavelength is then fixed by the material: silicon is opaque at the $1\,\mu$m wavelength of current detectors and of Cosmic Explorer Stage~1, so a silicon beamsplitter can only be used with the $2\,\mu$m laser of the Stage~2 configuration \citep[Table~1]{Reitze2019}; we therefore evaluate the dispersion at $2\,\mu$m. This is an inference from the given index, not a design fact: the Cosmic Explorer reference design lists a material only for the test masses (silica, or silicon for the $2\,\mu$m variant) and treats the beamsplitters as fused-silica optics \citep[Table~8.1 and Sec.~8.3.2]{Evans2021}, and the LIGO beamsplitter that the problem calls similar is fused silica \citep{GroteStadnik2019}, with $n\approx1.45$ at $1\,\mu$m; only its $6$~cm thickness matches LIGO's. The measured dispersion of silicon at $2.0\,\mu$m and $295$~K is $\mathrm{d}n/\mathrm{d}\lambda=-0.0382\,\mu\mathrm{m}^{-1}$ with $n=3.4535$ \citep{Frey2006}, giving $(\omega/n)\,\mathrm{d}n/\mathrm{d}\omega=0.0221$ and
\begin{equation}
\kappa_n=-0.0442 .
\end{equation}
This term has the same sign as $\kappa_l$ and contributes $4.4\%$ of the signal after the geometric weights of Step~4.\footnote{For the fused-silica beamsplitters of current detectors the index term is at the $10^{-3}$ level and is customarily dropped \citep{Vermeulen2021}; for silicon it is not. The temperature of the beamsplitter is specified neither by the problem nor, for the beamsplitter, by the Cosmic Explorer design \citep{Reitze2019}; the room-temperature dispersion is used, and the $123$~K value (interpolated from the same data, $\kappa_n=-0.0429$) would change the signal by $0.14\%$. This is recorded as an input uncertainty in the closing paragraph.}

\paragraph*{Step 4: Differential optical path of a Michelson interferometer with a thick beamsplitter}
Let the arms lie along $x$ and $y$, the beamsplitter center of mass at the origin, and the plate normal $\hat{\mathbf N}=(1,-1)/\sqrt2$ (incidence angle $\theta_i=45^\circ$, fixed by the geometry of the Michelson interferometer). The $50\%$ coating is on the face $\hat{\mathbf N}\cdot\mathbf r=-l/2$, which faces the laser and the $y$ arm; the opposite face is anti-reflection coated \citep[Fig.~1]{GroteStadnik2019}. Following both beams from the laser to the output port: the $x$ beam is transmitted through the substrate on the way out (one pass), and on return enters the substrate, reflects at the coating from the inside, and exits (two passes); the $y$ beam reflects at the coating on the way out and is transmitted through the substrate once on return. The substrate is therefore crossed three times by the $x$-arm beam and once by the $y$-arm beam, and the coating is used as a mirror from opposite sides by the two beams. Inside the plate, the ray is refracted by Snell's law to $\theta_r=\arcsin(\sin\theta_i/n)=11.66^\circ$ \citep[Sec.~1.5.1]{BornWolf1999}.

Rather than rely on any published closed form, the script traces both rays with exact vector Snell refraction and reflection, with the plate faces at $\pm l/2$ about the fixed center of mass and fixed reference wavefront planes at the laser, at the two end mirrors, and at the photodetector, and differentiates the round-trip optical path difference $\Delta\mathcal L=\mathcal L_x-\mathcal L_y$ numerically in 40-digit arithmetic.\footnote{For a plane wave the phase accumulated between two wavefront planes equals the optical path along any ray connecting them \citep[theorem of Malus and Dupin, Sec.~3.3.3]{BornWolf1999}, so the lateral displacement of the refracted beam does not enter the phase; the result is verified to be independent of the positions of the reference planes to machine precision, the exit rays are exactly antiparallel to the arms, and the calibration $\partial\Delta\mathcal L/\partial X_E-\partial\Delta\mathcal L/\partial Y_E=4$ for the end-mirror positions $X_E$, $Y_E$ is recovered. Reversing which face carries the coating exchanges the roles of the two arms and leaves $|\Delta\mathcal L|$ unchanged.} The numerical result recovers the closed form to better than $10^{-12}$:
%
\begin{equation}
\delta(\Delta\mathcal L)=2\left[n\cos\theta_r\,\delta l+l\,\frac{\partial(n\cos\theta_r)}{\partial n}\,\delta n\right],\qquad
\frac{\partial(n\cos\theta_r)}{\partial n}=\cos\theta_r+\frac{\sin^2\theta_i}{n^2\cos\theta_r},
\label{eq:geom}
\end{equation}
%
with $n\cos\theta_r=3.4278$ per encounter with the plate for the thickness term and $l\times1.0211$ per encounter for the index term. Each encounter contributes equally: on the outbound encounter the $x$ beam gains the plate-insertion path $l(n\cos\theta_r-\cos\theta_i)$ while the $y$ beam loses $l\cos\theta_i$ from the approaching coating; on the return encounter the $x$ beam gains the same insertion path plus $l\cos\theta_i$ from the receding coating while the $y$ beam gains the insertion path. The replacement of vacuum by glass ($-\cos\theta_i$) and the motion of the reflecting face ($+\cos\theta_i$) cancel, leaving $n\cos\theta_r$ per encounter.\footnote{In the no-refraction limit $\theta_r\to\theta_i$, Eq.~\eqref{eq:geom} reduces to $\sqrt2(n-\tfrac12)$ per encounter, which is the corrected beamsplitter expression of \citet[Eq.~5]{Vermeulen2021}; the earlier $\sqrt2\,n$ of \citet[Eq.~17]{GroteStadnik2019} omits the replacement of vacuum by glass. \citet{Vermeulen2021} then apply a refraction correction of about $6\%$ for their fused-silica beamsplitter; Eq.~\eqref{eq:geom} gives $6.1\%$ for $n=1.45$. For the $n=3.5$ beamsplitter of this problem the refraction correction is $19\%$, and the original $\sqrt2\,n$ overestimates the thickness coefficient by $44\%$ (Table~\ref{tab:results}).}

\paragraph*{Step 5: Finite light-travel time, direction dependence, and equivalent strain}
The two encounters with the plate are separated by the round-trip time $2L/c=267\,\mu$s, which at $200\,\mathrm{Hz}$ is a phase $2\omega L/c=0.335$~rad; the plate is hence not in the same state at the two encounters. Expressing the state of the plate at the first and second encounter as $\epsilon(t-2L/c)$ and $\epsilon(t)$, the signal at the photodetector is
\begin{equation}
\delta(\Delta\mathcal L)(t)=\mathcal A\left[\epsilon(t-2L/c)+\epsilon(t)\right],\qquad
\mathcal A=n\cos\theta_r\,\kappa_l\,l+l\,\frac{\partial(n\cos\theta_r)}{\partial n}\,\kappa_n\,n=-0.21515\,\mathrm{m},
\end{equation}
whose amplitude at $\omega$ is $2|\mathcal A|\cos(\omega L/c)\,\epsilon_0=0.42427\,\mathrm{m}\times\epsilon_0$, a $1.4\%$ reduction from the instantaneous value $2|\mathcal A|$.\footnote{The same finite-light-travel-time physics was found to matter for dark-photon searches with LIGO \citep{LVK2022}. Within each encounter the events of the two beams are separated by at most $1.4$~ns, a phase of $2\times10^{-6}$, so grouping them as ``early'' and ``late'' is exact at the precision required; the script uses the exact event times.}

The strain-noise curve is calibrated against gravitational waves, so the signal must be expressed through the same detector response function \citep{Rakhmanov2008}. For a Michelson interferometer with one round trip, a strain $h(t)=h_0\cos\omega t$ changes the round-trip optical path of the $x$ arm by $(c/2)\int_{t-2L/c}^{t}h\,\mathrm{d}t'$ and that of the $y$ arm by the opposite, so $|\Delta\mathcal L|=2Lh_0\,\mathrm{sinc}(\omega L/c)$ with $\mathrm{sinc}\,x=\sin x/x$, the well-known high-frequency correction to the long-wavelength response \citep[Eq.~14 at normal incidence]{Rakhmanov2008}. At $200\,\mathrm{Hz}$ and $L=40\,\mathrm{km}$, $\omega L/c=0.1677$ and $\mathrm{sinc}=0.99532$ (the script verifies this by direct quadrature). The equivalent strain amplitude of the dark-matter signal is therefore
\begin{equation}
h_0=\frac{2|\mathcal A|\cos(\omega L/c)}{2L\,\mathrm{sinc}(\omega L/c)}\,\epsilon_0\equiv\mathcal H\,\Lambda_\gamma^{-1},\qquad
\mathcal H=\frac{|\mathcal A|\cos(\omega L/c)}{L\,\mathrm{sinc}(\omega L/c)}\,\phi_0=2.1311\times10^{-4}\,\mathrm{GeV}.
\label{eq:H}
\end{equation}
Without the two light-travel-time factors $\mathcal H$ would be $2.1510\times10^{-4}\,\mathrm{GeV}$; their net effect is $-0.93\%$, a third-digit effect that the long arms of Cosmic Explorer make unavoidable.

Two further contributions depend on the direction of the dark-matter velocity, which is not known, and are therefore computed over the full sphere of directions (Fig.~\ref{fig:direction}). First, the end mirrors undergo the same thickness oscillation $\delta w/w=\kappa_l\,\delta\alpha/\alpha$ as the beamsplitter, and their reflecting faces move toward the beamsplitter by $\delta w/2$ with $w=40$~cm the Stage~2 test-mass thickness \citep[twice the aLIGO value,][]{Reitze2019,GroteStadnik2019}. The two arms cancel except for the phase difference $\mathbf k\cdot(\mathbf r_x-\mathbf r_y)$ of the field between the end mirrors at $\mathbf r_x$ and $\mathbf r_y$, which is at most $\sqrt2\,kL$ with $kL=1.29\times10^{-4}$ (reduced de~Broglie wavelength $3.1\times10^{8}$~m), and which enters in quadrature with the beamsplitter term. Second, the gradient of $\alpha$ exerts a force on every optic through the Coulomb-energy fraction $\kappa_M=1.7\times10^{-3}$ of the silicon atomic mass \citep{DamourDonoghue2010,GroteStadnik2019}, displacing each free mass by $\delta x=\kappa_M(v/\omega)\epsilon_0\sin(\omega t-\mathbf k\cdot\mathbf r)=0.31\,\mathrm{m}\times\epsilon_0$ along $\hat{\mathbf v}$. A common translation of all optics cancels in the differential output, but not exactly, because the optics are sampled at different times ($2L/c$ apart for the beamsplitter's two encounters). This residual reaches $3\%$ of the beamsplitter term for a wind along a plate normal, but it is $90^\circ$ out of phase and adds in quadrature. The full direction dependence of $h_0$ is $\pm3\times10^{-4}$ (relative), in the fourth significant digit, and the direction-averaged $\mathcal H$ is quoted above.

\begin{figure}[t]
\centering
\includegraphics[width=0.72\linewidth]{supporting/figs/direction_dependence.png}
\caption{Relative variation of the signal amplitude $h_0$ with the direction of the dark-matter velocity (azimuth measured from the $x$ arm, polar angle from the vertical), including the end-mirror thickness oscillation and the gradient-force displacement of all optics with their exact sampling times. The variation is at the $\pm3\times10^{-4}$ level and does not affect the third significant digit.}
\label{fig:direction}
\end{figure}

\paragraph*{Step 6: Detection statistics with finite coherence}
\emph{Monochromatic signal.} For $h(t)=h_0\cos(2\pi f_0t)$ on $0<t<T$ the Fourier amplitude near $f_0$ is $|\tilde h(f)|^2=(h_0^2T^2/4)\,\mathrm{sinc}^2[\pi(f-f_0)T]$, whose integral over $f$ is $h_0^2T/4$. With the one-sided noise convention, $\mathrm{SNR}^2=4\int_0^\infty|\tilde h|^2/S_h\,\mathrm{d}f=h_0^2T/S_h$, giving $\mathrm{SNR}=h_0\sqrt{T/S_h}$. The script also confirms this analytical expression by a Monte Carlo matched filter on white noise of known one-sided density ($0.713\pm0.018$ against the predicted $0.707$; a spurious factor $1/\sqrt2$ would have given $0.500$). Setting $\mathrm{SNR}=1$ defines the threshold $h_{\min}^{\rm coh}(T)=\sqrt{S_h/T}$.

\emph{Stochastic signal.} The dark-matter field is a superposition of waves with the velocity spread of the halo, so its frequency is spread by $\Delta\omega/\omega\sim(v/c)^2$ and its phase and amplitude wander after a coherence time. \citet{Derevianko2018} derived the field autocorrelation for a Maxwellian halo of velocity dispersion $\xi c$ seen from a detector moving through it at speed $v_g=\eta\,\xi c$:
\begin{equation}
\langle\phi(t)\phi(t+\tau)\rangle=\tfrac12\phi_0^2A(\tau)\cos[\omega'\tau+\psi(\tau)],\quad
A(\tau)=\frac{\exp\!\left[-\dfrac{\eta^2x^2}{2(1+x^2)}\right]}{(1+x^2)^{3/4}},\quad x=\frac{\tau}{\tau_c},\quad \tau_c=\frac{1}{\xi^2\omega},
\label{eq:corr}
\end{equation}
where $\tau$ is the time lag, $\tau_c$ the coherence time, $\omega'$ the Doppler-shifted frequency, and $\psi$ a slowly varying phase. The single velocity given in the problem is identified with the dispersion, $\xi c=230\,\mathrm{km\,s^{-1}}$, and the detector speed through the halo is taken equal to it ($\eta=1$, the fiducial case of \citealt{Derevianko2018}). Then $\tau_c=1352$~s.\footnote{Coherence-time conventions differ by factors of order $2\pi$: \citet{GroteStadnik2019} use $\tau_{\rm coh}=2\pi/(m_\phi v^2)=8495$~s for the same $v$, \citet{LVK2022} define the linewidth as $\tfrac12(v_0/c)^2f$. These are labels for the same physics as long as the detection statistic is computed consistently from Eq.~\eqref{eq:corr}, which is what is done here. A genuine, not merely notational, ambiguity is whether the problem's $230\,\mathrm{km\,s^{-1}}$ is the Maxwellian dispersion $\xi c$ or the streaming velocity $v_g$ of the detector through the halo. In the direct-detection literature a quoted ``dark-matter velocity'' of this size is the latter: the Earth's speed through the halo is $v_E\approx244\,\mathrm{km\,s^{-1}}$ in the classic parametrization of \citet[Eq.~3.6]{LewinSmith1996}, whose galactic rotation speed $v_0$ is precisely the problem's $230\,\mathrm{km\,s^{-1}}$, and $250\,\mathrm{km\,s^{-1}}$ with current values \citep{Baxter2021}, whereas the dispersion of the standard halo model is $v_0/\sqrt2=168\,\mathrm{km\,s^{-1}}$ for the recommended circular speed $v_0=238\,\mathrm{km\,s^{-1}}$ \citep{Baxter2021}. If the problem's $230\,\mathrm{km\,s^{-1}}$ is interpreted as the streaming velocity (the ``streaming reading'' below), the dispersion is not supplied by the problem and must be taken from the halo model, giving $\xi c=168\,\mathrm{km\,s^{-1}}$, $\eta=1.37$, $\tau_c=2525$~s, and a $0.7$~yr answer lower by $10\%$ (Table~\ref{tab:results}; the $1000$~s answer moves by $1.6\%$). The dispersion reading, in which the given value is the dispersion and $\eta=1$ so that both velocities equal the given $230\,\mathrm{km\,s^{-1}}$, is adopted as primary because it is internally consistent with the problem statement and uses only its own number (it is not a dynamically self-consistent halo: in an equilibrium, non-rotating halo the Jeans equation ties the dispersion to the circular speed, $\sigma_{1\mathrm D}\approx v_c/\sqrt2$, so a streaming speed equal to the dispersion at $230\,\mathrm{km\,s^{-1}}$ would require a halo about $40\%$ hotter than the Milky Way's rotation curve allows, a net halo rotation of order $70\,\mathrm{km\,s^{-1}}$, or a local halo out of equilibrium, such as unrelaxed merger debris, for which the Jeans relation does not apply). Note that $230\,\mathrm{km\,s^{-1}}$ is a standard value for the streaming speed and for $v_0$, but about $40\%$ above the one-dimensional dispersion of the standard halo model, which is why the streaming reading, the literature-faithful alternative reported alongside, has the longer coherence time. Both readings assume a dispersion-supported halo. A rotation-supported component from dissipative dark-matter models would have far smaller velocities in both roles and a correspondingly narrower line.} At $\tau=1000$~s, $A=0.60$: the field is already substantially decorrelated over the shorter of the two observation times, so the sharp division into a ``coherent'' and an ``incoherent'' regime is only an approximation there.

For a Gaussian stochastic signal the optimal statistic is quadratic, and its sensitivity to the coupling is $\sigma_\gamma\propto\left[\sum_p\left(\langle|\tilde\phi_p|^2\rangle/\tilde\rho_p\right)^2\right]^{-1/4}$ \citep[Eq.~15]{Derevianko2018}, where $\langle|\tilde\phi_p|^2\rangle$ is the expected power of the field in the $p$-th frequency bin of the discrete Fourier transform and $\tilde\rho_p$ the noise power in the same bin (white noise, so independent of $p$). The expected power spectrum of a record of duration $T$ is the Fourier transform of the correlation function weighted by the triangular window $(1-|\tau|/T)$, so by Parseval the sum is $\propto\int_0^T(1-\tau/T)^2A(\tau)^2\,\mathrm{d}\tau$, and the phase $\psi$ drops out. Normalizing to the monochromatic threshold, the finite-coherence threshold is
\begin{equation}
h_{\min}(T)=\sqrt{\frac{S_h}{T}}\;R(T),\qquad
R(T)=\left[\frac{T/3}{\int_0^T(1-\tau/T)^2A(\tau)^2\,\mathrm{d}\tau}\right]^{1/4},
\label{eq:R}
\end{equation}
which tends to $1$ for $T\ll\tau_c$ (a fully coherent signal, for which the monochromatic threshold applies) and to $[T/(3I\tau_c)]^{1/4}$ with $I=\int_0^\infty A^2\,\mathrm{d}x=0.747$ for $T\gg\tau_c$, reproducing the familiar $\propto(T\tau_c)^{-1/4}$ law with its constant derived rather than assumed.\footnote{Eq.~\eqref{eq:R} uses the window-averaged power spectrum as the coherent reference. A search that scans the template frequency finely recovers the on-bin value in the coherent limit, which would raise the incoherent-regime constant by $(3/2)^{1/4}=1.107$, the factor $3/2$ being the ratio of the on-bin to the bin-averaged value of $\sum_p\langle|\tilde h_p|^2\rangle^2$ for a monochromatic signal (bin-averaged $\sum_p\mathrm{sinc}^4=\pi^{-1}\!\int\mathrm{sinc}^4x\,\mathrm{d}x=2/3$, against $1$ on bin). We also note that evaluating \citet{Derevianko2018}'s Eq.~(15) with his own correlation function, as done here, gives an asymptotic constant smaller than the $2.4$ quoted in his Eq.~(16), which we could not reconcile with the stated normalization of his line-shape function. Both are $\mathcal O(1)$ choices inherent to defining an SNR for a stochastic signal and are absent from the problem statement. The two-regime idealization the problem implies is reported alongside in Table~\ref{tab:results}.} Numerically $R(1000\,\mathrm{s})=1.029$ and $R(0.7\,\mathrm{yr})=9.241$. The two-regime idealization with $\tau_{\rm coh}=2\pi/(m_\phi v^2)$ would give $1$ and $(T/\tau_{\rm coh})^{1/4}=7.14$ (Fig.~\ref{fig:R}).

\begin{figure}[t]
\centering
\includegraphics[width=0.66\linewidth]{supporting/figs/R_of_T.png}
\caption{Threshold factor $R(T)=h_{\min}(T)/\sqrt{S_h/T}$ from Eq.~\eqref{eq:R} for the two readings of the given velocity, compared with the sharp two-regime law implied by the problem. $R=1$ is the fully coherent limit $T\ll\tau_c$, in which the threshold is the monochromatic $\sqrt{S_h/T}$. The curves depart from it once $T$ approaches $\tau_c$. Vertical lines mark $T=1000$~s and $0.7$~yr.}
\label{fig:R}
\end{figure}

\paragraph*{Step 7: The coupling detection threshold}
Combining Eqs.~\eqref{eq:H} and \eqref{eq:R} at $\mathrm{SNR}=1$,
\begin{equation}
\Lambda_\gamma^{-1}\big|_{\min}=\frac{h_{\min}(T)}{\mathcal H}=\frac{\sqrt{S_h/T}\,R(T)}{\mathcal H}.
\label{eq:final}
\end{equation}
With $\sqrt{S_h/T}=6.3246\times10^{-27}$ ($T=1000$~s) and $4.2553\times10^{-29}$ ($T=0.7$~yr):
\begin{equation}
\Lambda_\gamma^{-1}\big|_{\min}=3.05\times10^{-23}\,\mathrm{GeV}^{-1}\ (1000\,\mathrm{s}),\qquad
\Lambda_\gamma^{-1}\big|_{\min}=1.85\times10^{-24}\,\mathrm{GeV}^{-1}\ (0.7\,\mathrm{yr}).
\end{equation}
Table~\ref{tab:results} shows how the forecast detection thresholds vary under the alternative setups and assumptions discussed above. The largest single effect is the beamsplitter geometry (Step~4), which lowers the thickness coefficient by $31\%$ relative to the uncorrected literature formula; together with the index term and the light-travel-time factors the signal is $28\%$ lower, and the threshold coupling correspondingly higher by a factor $1.39$. The finite-coherence statistics (Step~6) are the next largest for the $0.7$~yr observation ($29\%$ between the two-regime idealization and Eq.~\eqref{eq:R}).

\begin{table}[t]
\centering
\caption{Smallest detectable $\Lambda_\gamma^{-1}$ (GeV$^{-1}$) at SNR~$=1$ and $200$~Hz. ``Exact'' is Eqs.~\eqref{eq:geom}, \eqref{eq:H}, and \eqref{eq:R}. The last two rows use the published beamsplitter expressions with the two-regime idealization. Systematics on the first row: relativistic correction to the Bohr scaling, $-0.8\%$ to $-1.5\%$; direction of the dark-matter velocity, $\pm0.03\%$; beamsplitter temperature, $+0.14\%$.}
\label{tab:results}
\small
\begin{adjustbox}{max width=\linewidth}
\begin{tabular}{llcc}
\toprule
signal model & time-domain model & $T=1000$~s & $T=0.7$~yr \\
\midrule
exact (this work) & Eq.~\eqref{eq:R}, $\xi c=v_g=230$~km/s & $3.05\times10^{-23}$ & $1.85\times10^{-24}$ \\
exact & Eq.~\eqref{eq:R}, $v_g=230$, $\xi c=168$~km/s (SHM) & $3.00\times10^{-23}$ & $1.66\times10^{-24}$ \\
exact & two regimes, $\tau_{\rm coh}=2\pi/(m_\phi v^2)$ & $2.97\times10^{-23}$ & $1.43\times10^{-24}$ \\
\citet[Eq.~5]{Vermeulen2021}, no refraction & two regimes, as above & $2.49\times10^{-23}$ & $1.19\times10^{-24}$ \\
\citet[Eq.~18]{GroteStadnik2019} & two regimes, as above & $2.13\times10^{-23}$ & $1.02\times10^{-24}$ \\
\bottomrule
\end{tabular}
\end{adjustbox}
\end{table}

\paragraph*{What the problem idealizes}
Three physical features of Cosmic Explorer are outside the problem as posed and are therefore not in the primary answer. (i) \emph{Fabry-Perot arms.} The strain curve of \citet{Reitze2019} presupposes arm cavities (its $1.4$ to $2$~MW circulating power exists only inside them), whereas a beamsplitter signal is not resonantly enhanced by arm cavities. The effective signal is suppressed by $1/N_{\rm eff}$, with $N_{\rm eff}$ the effective number of round trips in the arm cavities, of order $10^{-2}$ for a LIGO-like finesse at low frequency and rising linearly with frequency above the cavity bandwidth \citep[Eqs.~20 and 21]{GroteStadnik2019}. The problem gives no finesse or signal-recycling parameters, so it can only intend the response function of a simple Michelson interferometer, which is what was used. In a Fabry-Perot Michelson interferometer the end-mirror term of Step~5, which is resonantly enhanced, would also compete with the beamsplitter term \citep{GroteStadnik2019}. (ii) \emph{Annual modulation.} Over $0.7$~yr the Earth's orbital velocity ($\pm30\,\mathrm{km\,s^{-1}}$) modulates both the line center and, through $v^2$, the coherence time by tens of percent \citep{Vermeulen2021,Baxter2021}. The problem fixes a single $v$. (iii) \emph{The overdensity scenario} (footnote to Step~1). None of these change the calculation requested; they change what ``Cosmic Explorer'' would actually measure.

\paragraph*{Is the problem well posed?}
As a statement it is: every input is specified and the requested quantity is defined. At three significant digits it is underdetermined, for reasons that are physical rather than notational. The beamsplitter model (refraction, index variation, the two-sided use of the coating) moves the answer by tens of percent between defensible published treatments; the relativistic correction to the size scaling of silicon and the scheme-dependent radiative shift of $m_e$ are each $1\%$ effects that no closed-form input can fix; and the meaning of ``SNR~$=1$ with an observation time'' for a partially coherent stochastic signal carries $\mathcal O(1)$ constants that the statement does not fix. The statement is also not closed in its inputs: the line shape of a wave-like dark-matter signal depends on two velocities, the dispersion of the halo and the streaming speed of the detector through it, and the problem supplies one number without saying which it is. Every way of filling the second slot is an assumption external to the problem, and the choices span $10\%$ in the $0.7$~yr answer (Table~\ref{tab:results}). The primary answer resolves this by setting both velocities equal to the given value, the only option that adds no information, at the price of a velocity pair that no equilibrium model of the local halo produces (Step~6 footnote). The beamsplitter material and laser wavelength are likewise not given, although the index term ($4.4\%$ of the signal) and the relativistic correction ($1\%$) require them at the three significant digits requested, and already at two. They are inferred from $n=3.5$ (silicon at $2\,\mu$m), an index that matches neither the LIGO beamsplitter the statement calls similar ($n=1.45$) nor the fused-silica beamsplitters of the Cosmic Explorer design \citep{Evans2021}. The beamsplitter temperature ($0.14\%$) is also unspecified. The exact chain above removes every approximation that can be removed and states each remaining choice explicitly. Agreement with a reference answer beyond the second digit would test which conventions the reference adopted, not the physics.

\paragraph*{Filled \texttt{template.py}}
\lstinputlisting[style=answerstyle,caption={\texttt{template.py}}]{template.py}

\bibliographystyle{unsrtnat}
\bibliography{supporting/refs}
